\documentclass[11pt]{scrartcl}

\usepackage[stixtwo]{fontsetup}
\usepackage{amsmath}

\begin{document}
Teljes indukcióval megmutatjuk, hogy $1 \le d_n \le 4$ minden
$n \in \mathbb{N}^+$ esetén. Mivel $d_1 = 2$, ezért $n = 1$-re az állítás
teljesül. Ha valamilyen $n$-re teljesül, azaz $1 \le d_n \le 4$, akkor
$\frac{1}{4} \le \frac{1}{d_n} \le 1$, így $1 \le \frac{4}{d_n} \le 4$, és
ezért $1 \le 5 - \frac{4}{d_n} = d_{n + 1} \le 4$.

Teljes indukcióval megmutatjuk, hogy $d_n$ monoton növő, azaz
$d_n \le d_{n + 1}$ minden $n \in \mathbb{N}^+$ esetén. Mivel $d_2 = 3$,
ezért $n = 1$-re az állítás teljesül. Ha valamilyen $n$-re teljesül, azaz
$d_n \le d_{n + 1}$, akkor mivel $d_n$ pozitív, ezért
$\frac{1}{d_n} \ge \frac{1}{d_{n + 1}}$, és így
\[
	d_{n + 1} = 5 - \frac{4}{d_n} \le 5 - \frac{4}{d_{n + 1}} = d_{n + 2},
\]
vagyis $n + 1$-re is teljesül az egyenlőtlenség.

Mivel a sorozat korlátos és monoton, ezért konvergens, határértéke a
rekurzió valamelyik fixpontja. Az $A = 5 - \frac{4}{A}$ egyenlet megoldásai
$A = 1$ és $A = 4$. Mivel $d_n$ monoton növő és $d_2 = 3$, a határérték nem
lehet $1$. Tehát $\lim d_n = 4$.
\end{document}
